\section{Introduction}

Borrowing words from one language (or the source language \( L_2 \)) to another
(the borrowing language \( L_1 \)) is a common linguistic phenomenon. This is
more common in locales where multiple languages are being used simultaneously.

Consider the contemporary Filipino as an example, which routinely borrows
lexical items from English for use in the education, medicine, and
administrative sectors. Consider the specific example of going to school: one
might use ``paaralan'' which is the native Filipino word for ``school''. Or
borrow by using ``eskuwelahan'' from the Spanish word ``escuela'', or ``iskul''
or ``skul' from the English word ``school''.

For the most part, adapting borrowed words is a phonetic task: ``how do I say
this word, given the phonology of the borrowing language?''. An example of this
in Filipino, adapting the word ``standby''. First, Filipino consinant clusters,
like \emph{st}, are not natural. Second, the transition from the /m/ to the /b/
sound is more natural than /n/ to /b/. Finally, Filipino has a strict set of
native vowel, hence the \emph{y} which sounds like /\textipa{aI}/ is written as
\emph{ay} which sounds like /\textipa{ay}/. All of which contribute to the
adaptation, which yields with ``istambay''.

However, adapting loanwords is not purely phonetic, since how the word is also
spelt, or its orthography, also influences adaptation. Consider the English word
``biscuit'' /\textipa{'bIs.kIt}/ whose adapted form is ``biskwit'' /\textipa{/'bis.kwIt/}/.
This example shows that silent letters have a tendency to be realized
when adapting, and the reference of this adaptation is through the orthography.

Add more literature review here or whatever

Finally, even though Filipino has a orthography manual publicly available, no
software solution is yet to be implemented for writing adapted loanwords.

